% Hardware, software, and the two computers of the Day 7 lab.
\begin{frame}{Hardware and software}
  \begin{columns}[T]
    \begin{column}{0.46\textwidth}
      \textbf{Hardware for each group}
      \begin{itemize}
        \item Raspberry Pi 5 (8 GB) with the active cooler and the card of
              your group: \code{pi-NN}.
        \item Power supply, 27 W.
        \item Camera Module 3 with its cable.
        \item Laptop in the Wi-Fi \code{edgeai-lab}.
      \end{itemize}
    \end{column}
    \begin{column}{0.50\textwidth}
      \textbf{Software}
      \begin{itemize}
        \item On the board: the environment \code{\textasciitilde/tflite\_env}
              with LiteRT and ONNX Runtime.
        \item On the laptop: a new environment with the file
              \code{requirements.txt}, for the notebook.
        \item Netron: \code{netron.app} in a browser.
      \end{itemize}
    \end{column}
  \end{columns}
  \vskip 0.6em
  \small
  \renewcommand{\arraystretch}{1.15}
  \begin{tabular}{@{}L{0.20\textwidth}L{0.72\textwidth}@{}}
    Lab files & \code{Labs/day07/} on the laptop,
                \code{\textasciitilde/edgeai/day07/} on the board \\
    Each new terminal on the board &
                \code{cd \textasciitilde/edgeai/day07} and
                \code{source \textasciitilde/tflite\_env/bin/activate} \\
    The camera cable & Do not connect or remove it when the power is on. \\
  \end{tabular}
\end{frame}

\begin{frame}{Two computers, one model}
  \begingroup\centering
  \begin{tikzpicture}[font=\footnotesize, >={Stealth[length=5pt]},
      b/.style={draw=coolgray, rounded corners=3pt, fill=boxgray,
                minimum width=1.75cm, minimum height=1.0cm, align=center},
      lab/.style={font=\scriptsize, text=coolgray, align=center}]
    \node[b] (a) at (0,0)    {PyTorch\\model};
    \node[b, draw=darkorange] (b) at (2.15,0) {Export\\(Part C)};
    \node[b] (c) at (4.3,0)  {Model\\files};
    \node[b, draw=darkcyan] (d) at (6.45,0) {Runtime\\(Parts B, D)};
    \node[b, draw=cardinalred] (e) at (8.6,0)  {Latency\\table};
    \foreach \a/\b in {a/b, b/c, c/d, d/e} { \draw[->] (\a) -- (\b); }
    \node[lab] at (0,-0.95)    {MobileNetV2,\\141 layers};
    \node[lab] at (2.15,-0.95) {static, dynamic,\\LiteRT};
    \node[lab] at (4.3,-0.95)  {\texttt{.onnx}, \texttt{.tflite},\\copy with \texttt{scp}};
    \node[lab] at (6.45,-0.95) {LiteRT,\\ONNX Runtime};
    \node[lab] at (8.6,-0.95)  {threads,\\precision};
    \draw[coolgray, dashed] (-1.0,0.75) rectangle (5.32,-1.45);
    \draw[coolgray, dashed] (5.43,0.75) rectangle (9.6,-1.45);
    \node[lab] at (2.15,1.0) {the laptop};
    \node[lab] at (7.5,1.0)  {the Raspberry Pi};
  \end{tikzpicture}\par\endgroup
  \vskip 0.4em
  \small
  \begin{itemize}
    \item The export needs PyTorch and the converters. They stay on the
          laptop. The board needs only the two runtimes.
    \item The folder \code{models/} already has three files: the model of
          the kit lab in \code{uint8}, and the model of Part C in
          \code{int8} for each runtime.
    \item You make the \code{float32} files in Part C and copy them to the
          board.
  \end{itemize}
  \keypoint{This is the path of each deployment: train and export on a
    large computer, run a small file on the device.}
\end{frame}
